\begin{solapartat}
\begin{minipage}[t]{0.68\linewidth}
$E = k\dfrac{q}{r^2}$; \ $\vec{E}_x = (4000\ \mathrm{N/C})\,\vec{\imath}$

A l'eix vertical el camp elèctric s'anul·la per simetria.

\medskip
\punts{0,4} $E_x = 2k\dfrac{Q}{l^2}\cos\theta \rightarrow Q = \dfrac{E_x l^2}{2k\cos\theta}$

\medskip
\punts{0,2} $\cos\theta = \dfrac{4}{\sqrt{4^2 + 2^2}} = \dfrac{4}{\sqrt{20}}$
\end{minipage}\hfill
\begin{minipage}[t]{0.28\linewidth}
\vspace{0pt}
\figura[1]{sol_fig1.png}
\end{minipage}

\medskip
\punts{0,2} $Q = \dfrac{4\times 10^{3}\cdot 20}{2\cdot 8,99\times 10^{9}\cdot\dfrac{4}{\sqrt{20}}} = 4,97\times 10^{-6}\ \mathrm{C} = 4,97\ \mu\mathrm{C}$ i és negativa \punts{0,2} perquè el camp elèctric total va dirigit en el sentit positiu de l'eix X.
\end{solapartat}

\begin{solapartat}
\punts{0,2} $F_y = 0$ per simetria.

\medskip
\punts{0,4}
$\left.\begin{aligned}
&F_x = q_p E_x\\
&E_x = 2k\frac{Q}{l^2}\cos\theta\\
&Q = 2\times 10^{-6}\ \mathrm{C}
\end{aligned}\right\} \rightarrow F_x = 2k\dfrac{Qq_p}{l^2}\cos\theta$

\medskip
$F_x = 2\cdot 8,99\times 10^{9}\cdot\dfrac{2\times 10^{-6}\cdot 1,6\times 10^{-19}}{20}\cdot\dfrac{4}{\sqrt{20}} = 2,57\times 10^{-16}\ \mathrm{N}$

\medskip
En aquest cas, la càrrega $Q$ (protó) és positiva, el camp elèctric total va dirigit en el sentit negatiu de l'eix $X$ i la força sobre un protó també.

\medskip
\punts{0,4} $\vec{F} = \left(-2,57\times 10^{-16}\ \mathrm{N}\right)\vec{\imath}$
\end{solapartat}
